% !TeX root = ../main.tex
\chapter{研究目标与研究内容}

\section{研究目标}
给出本课题拟解决的核心问题，并将总体目标分解为可验证的具体目标。
评价指标应明确测量对象、基线、测试条件与计算方法；指标目标值需要依据调研和初步实验确定。

\section{主要研究内容}
\subsection{研究内容一：问题建模与分析}
说明研究对象如何表示、有哪些变量和约束，以及拟采用的分析方法。
写清本项工作的输入、输出，以及它与后续工作的依赖关系。

\subsection{研究内容二：方法或系统设计}
说明拟研究的方法、关键模块及其作用，指出与现有工作的差异。
为每个设计选择说明依据，并列出需要实验检验的假设。

\subsection{研究内容三：实现与验证}
说明拟完成的原型、数据或实验平台，以及验证范围。
区分正确性验证、性能比较和消融实验，保证评价能够回答研究问题。

\section{拟解决的关键问题与预期创新点}
将关键问题与研究内容一一对应。预期创新点应具体到新的问题表述、方法机制或验证结论，
并说明如何通过对比实验或分析验证；不要将一般性的工程实现直接等同于创新。
